% Generated draft: requires mathematical and source review.
\documentclass[11pt]{article}
\usepackage[margin=25mm]{geometry}
\usepackage{fontspec}
\setmainfont{Latin Modern Roman}
\usepackage{amsmath,amssymb,mathtools,unicode-math}
\setmathfont{Latin Modern Math}
\usepackage{enumitem,needspace,xurl,hyperref}
\setlength{\parindent}{0pt}
\setlength{\parskip}{5pt}
\setlength{\emergencystretch}{3em}
\newcommand{\N}{\mathbb{N}}
\newcommand{\Z}{\mathbb{Z}}
\newcommand{\Q}{\mathbb{Q}}
\newcommand{\R}{\mathbb{R}}
\newcommand{\C}{\mathbb{C}}
\newcommand{\F}{\mathbb{F}}
\newcommand{\A}{\mathbb{A}}
\newcommand{\PP}{\mathbb{P}}
\newcommand{\E}{\mathbb{E}}
\newcommand{\eps}{\varepsilon}
\newcommand{\dd}{\,\mathrm{d}}
\begin{document}
\section*{UnsolvedMath ID: 30006492}
\textbf{Source problem code:} catalogue entry
\subsection{Definitions and mathematical statement}
Let \(X\) be a set. A set-theoretic solution of the Yang--Baxter equation is a map \(t\colon X\times X\to X\times X\) satisfying
\[
(t\times\operatorname{id}_X)(\operatorname{id}_X\times t)(t\times\operatorname{id}_X)
=
(\operatorname{id}_X\times t)(t\times\operatorname{id}_X)(\operatorname{id}_X\times t)
\quad\text{on }X^3.
\]
This equation alone does not specify any additional bijectivity or nondegeneracy assumptions required in the cited report.

The supplied problem concerns a \emph{welded pair} \((X,r,s)\), its induced representations of the welded braid group \(WB_n\) on \(n\) strands, and a \emph{derived solution} \(r'\) associated with \(r\). The supplied extract does not give the defining compatibility conditions for a welded pair, the construction of \(r'\), or the induced representation formulas. These are unresolved prerequisites for a fully specified statement.

\textbf{Questions retained from the supplied record.}
\begin{enumerate}
\item When the representations induced by \((X,r,s)\) are enriched by weights computed using \(2\)-cocycles, can the second component \(s\) affect their isomorphism classes? Here representation-theoretic relevance is interpreted as dependence of the isomorphism class on \(s\); the source does not specify the cocycle theory, coefficient objects, admissible weights, or how weights are to be compared when \(s\) changes.
\item Can the representations under consideration be replaced, up to isomorphism, by representations induced by a welded pair \((X,r',q)\) for some second component \(q\)? Writing \(\omega\) and \(\widetilde\omega\) only as placeholders for unspecified admissible weight data, the intended comparison has the schematic form
\[
\rho_n^{r,s;\omega}\cong\rho_n^{r',q;\widetilde\omega},
\qquad (X,r',q)\text{ a welded pair}.
\]
This display does not define the representations or assert that such weight data exist. For representations on modules \(V\) and \(V'\) over a common coefficient ring, an isomorphism would mean an invertible linear map \(T\colon V\to V'\) satisfying
\[
T\rho_n^{r,s;\omega}(g)
=\rho_n^{r',q;\widetilde\omega}(g)T
\qquad\text{for every }g\in WB_n.
\]
The coefficient ring, modules, and permitted changes of cocycles are not supplied. The original wording of this second question does not explicitly mention weights, whereas the cleaned statement connects it to the weighted representations. Both readings must remain open pending examination of the surrounding source.
\end{enumerate}
The source does not state the admissible class of \(X,r,s\), the range of \(n\), or whether a single \(q\) and compatible weight data must work simultaneously for all strand numbers. Thus the quantifiers needed to turn these questions into a precise universal assertion remain unresolved.

The trailing phrase ``Can we weld a YBE solution?'' may be a contribution title or an additional question. If it is a separate mathematical question, it asks about extending a Yang--Baxter solution to a welded pair, but its hypotheses and intended scope are not supplied. It is retained here without asserting an unrestricted extension conjecture.
\subsection{Short English statement}
The problem asks whether cocycle weights make the second component of a welded pair affect its braid-group representations, and whether those representations can still be obtained using the derived first component. The supplied extract omits the constructions, weight conditions, and several quantifiers needed for a complete formulation. Its open status is supplied dated curation, not an independently verified conclusion.
\subsection{Sources}
\begin{itemize}
\item Supplied local UnsolvedMath registry record, permanent ID \(30006492\), source problem code \texttt{catalogue entry}, title \emph{Cocycle-Weighted Representations of Welded Braid Groups}. The supplied collection is \emph{Oberwolfach Reports Open Problems}.
\item Supplied source citation: \emph{Mini-Workshop: The Yang--Baxter Equation and Representations of Braid Groups (2026)}, Oberwolfach Reports, pp.\ 2633--2672, DOI \url{https://doi.org/10.4171/owr/2025/49}. The extract includes page number 2655. The citation's parenthetical year 2026 differs from the report designation 49/2025 and the supplied proposed year 2025; this discrepancy is unresolved.
\item The extracted passage names Victoria Lebed and identifies joint work with Ilaria Colazzo, Jo\~ao Faria Martins, and Senne Trappeniers. The registry does not separately identify a proposer. The phrase \emph{Can we weld a YBE solution?} appears alongside this attribution, but its role is unclear from the extract.
\item Supplied literature source: \emph{Oberwolfach Report 49/2025}, \url{https://doi.org/10.4171/owr/2025/49}.
\item Supplied literature source: \emph{Yang--Baxter 2-cocycle invariants}, \url{https://doi.org/10.18910/73356}. The extract refers to references [1] and [4] without their bibliography entries; their identities are not established by the supplied record.
\item Literature assessment checked on \texttt{2026-08-22}: this is supplied dated curation, not independently verified. It labels the questions open and reports that the questions about weighted representations and replacement of the second welded component follow an unweighted reduction theorem, with no general resolution found in that supplied assessment. The background expressly requests independent verification of time-sensitive status and cited claims.
\item The registry labels the statement \texttt{corrected\_verified} and says that the official report was checked and notation repaired. These are supplied verification claims; no independent source consultation or mathematical verification is asserted in this draft.
\end{itemize}
\end{document}
